%% ============================================================
%%  The Course List of Facts.
%%  Back-matter reference page. The facts use stable keys and
%%  auto-generated F## identifiers supplied by \coursefact.
%% ============================================================
\documentclass[main.tex]{subfiles}
\begin{document}

\clearpage
\phantomsection\addcontentsline{toc}{chapter}{The Course List of Facts}
\markboth{The Course List of Facts}{The Course List of Facts}
\thispagestyle{plain}
{\LARGE\bfseries\hyphenpenalty=10000\exhyphenpenalty=10000
 The Course List of Facts\par}\nobreak\smallskip
\noindent\rule{\linewidth}{0.4pt}\par\medskip
\noindent These are the results the instructor lists as quotable on homework and exams. The list is
not exhaustive. Each fact has a stable key and an automatically generated identifier, and each carries
a pointer to where it lives in our notes. ``Quotable'' means you may use it without reproving it, but
you should still know \emph{why} it holds.

\par\medskip\noindent{\large\bfseries Number Systems}\par\nobreak\smallskip
\begin{itemize}[leftmargin=1.4em,itemsep=2pt,topsep=2pt]
  \coursefact{q-ordered-field}
    {$\mathbb Q$ is an ordered field: all the field and order rules, and valid manipulations of
     (in)equalities, may be used without proof.}
    {\xrefL{1.3.4}}
  \coursefact{q-no-sup}
    {There exist nonempty subsets of $\mathbb Q$ that are bounded above but have no supremum in
     $\mathbb Q$.}
    {\xrefL{1.4.2}}
  \coursefact{lub}
    {$\mathbb R$ is an ordered field containing $\mathbb Q$ as a subfield and satisfying the
     least-upper-bound property: every nonempty bounded-above subset has a supremum in $\mathbb R$.}
    {\xrefL{1.6.1}}
  \coursefact{archimedean}
    {$\mathbb R$ is Archimedean: if $x>0$ then $nx>y$ for some $n\in\mathbb N$; equivalently
     $\inf\{1/n\}=0$.}
    {\xrefL{1.7.1}}
  \coursefact{q-dense}
    {$\mathbb Q$ is dense in $\mathbb R$: every real interval contains a rational.}
    {\xrefL{1.7.2}}
\end{itemize}

\par\medskip\noindent{\large\bfseries Countable and Uncountable Sets}\par\nobreak\smallskip
\begin{itemize}[leftmargin=1.4em,itemsep=2pt,topsep=2pt]
  \coursefact{nzq-countable}
    {$\mathbb N$, $\mathbb Z$, and $\mathbb Q$ are countable.}
    {\xrefL{1.7.2}; \xrefL{2.1.4}, \xrefL{2.1.7}}
  \coursefact{subset-countable}
    {A subset of a countable set is finite or countable.}
    {\xrefL{2.1.9}}
  \coursefact{countable-union}
    {A countable union of countable sets is countable.}
    {used in the notes; not stated as a separate object}
  \coursefact{finite-product}
    {A finite product of countable sets is countable.}
    {used in the notes; not a separate object}
  \coursefact{binary-uncountable}
    {The set of all sequences of $0$'s and $1$'s is uncountable.}
    {\xrefL{2.1.10}}
  \coursefact{r-uncountable}
    {$\mathbb R$ is uncountable.}
    {\xrefL{2.1.13}}
\end{itemize}

\par\medskip\noindent{\large\bfseries Open and Closed Subsets}\par\nobreak\smallskip
\begin{itemize}[leftmargin=1.4em,itemsep=2pt,topsep=2pt]
  \coursefact{balls-open}
    {The open balls of $(X,d)$ are open.}
    {\xrefL{2.3.5}}
  \coursefact{cluster-infinite}
    {If $p$ is a cluster point of $E$, every ball around $p$ contains infinitely many points of $E$.}
    {\xrefL{3.2.1}}
  \coursefact{open-iff-closed}
    {A subset is open in $(X,d)$ iff its complement is closed.}
    {\xrefL{3.2.8}}
  \coursefact{union-open}
    {Any union of open sets is open; a finite intersection of open sets is open.}
    {\xrefL{3.1.2}}
  \coursefact{inter-closed}
    {Any intersection of closed sets is closed; a finite union of closed sets is closed.}
    {\xrefL{3.2.10}}
  \coursefact{closure-closed}
    {If $E'$ is the set of cluster points of $E$, then $E\cup E'$ is closed.}
    {\xrefL{4.4.1}}
\end{itemize}

\par\medskip\noindent{\large\bfseries Compactness}\par\nobreak\smallskip
\begin{itemize}[leftmargin=1.4em,itemsep=2pt,topsep=2pt]
  \coursefact{compact-absolute}
    {Compactness is absolute: $K$ is compact relative to $X$ iff relative to any $Y\supseteq X$; one
     may simply say $(K,d)$ is compact.}
    {\xrefL{3.4.4}}
  \coursefact{compact-closed}
    {Compact subsets of metric spaces are closed and bounded.}
    {\xrefL{4.1.1}}
  \coursefact{closed-in-compact}
    {Closed subsets of compact spaces are compact.}
    {\xrefL{4.1.3}}
  \coursefact{nested-compact}
    {The intersection of nonempty, nested compact sets is nonempty.}
    {\xrefL{4.2.2}}
  \coursefact{infinite-cluster}
    {An infinite subset of a compact space has a cluster point.}
    {\xrefL{4.3.3}}
  \coursefact{heine-borel}
    {A subset of $\mathbb R^k$ is compact iff it is closed and bounded.}
    {\xrefL{4.3.5}}
\end{itemize}

\par\medskip\noindent{\large\bfseries Sequences}\par\nobreak\smallskip
\begin{itemize}[leftmargin=1.4em,itemsep=2pt,topsep=2pt]
  \coursefact{limit-tail}
    {If $x_n\to p$, every ball around $p$ contains all but finitely many $x_n$.}
    {\xrefL{5.1.3}}
  \coursefact{limit-unique}
    {A sequence in a metric space has at most one limit.}
    {\xrefL{5.1.3}}
  \coursefact{conv-bounded}
    {A convergent sequence is bounded.}
    {\xrefL{5.1.3}}
  \coursefact{cluster-seq}
    {If $p$ is a cluster point of $E$, some sequence in $E$ converges to $p$.}
    {\xrefL{5.1.3}}
  \coursefact{limit-algebra}
    {Limits in $\mathbb R^k$ respect the algebraic operations.}
    {\xrefL{5.2.1}}
  \coursefact{coordinatewise}
    {A sequence in $\mathbb R^k$ converges iff each coordinate converges.}
    {\xrefL{5.2.2}}
  \coursefact{compact-subseq}
    {Every sequence in a compact metric space has a convergent subsequence; every bounded sequence in
     $\mathbb R^k$ does too.}
    {\xrefL{5.3.3}}
  \coursefact{conv-cauchy}
    {Every convergent sequence is Cauchy; every Cauchy sequence is bounded; if a subsequence of a
     Cauchy sequence converges to $p$, so does the sequence.}
    {\xrefL{5.4.2}, \xrefL{5.5.3}}
  \coursefact{complete}
    {Every Cauchy sequence in $\mathbb R^k$ or in a compact metric space converges: these are
     complete.}
    {\xrefL{5.5.3}}
  \coursefact{monotone-conv}
    {A bounded, monotonic sequence in $\mathbb R$ converges.}
    {\xrefL{6.2.2}}
\end{itemize}

\end{document}
